%%% FIELD proof-uniform-normal-bound
Fix \((G_n,Y_n)\in\mathcal M_n(d;B,c)\), and put
\[
 s_n^2=n^2\sigma_n^2.
\]
Membership in \(\mathcal M_n(d;B,c)\), specifically \(\mathrm{ass:variance\text{-}floor}\), gives
\[
 s_n^2\ge c n\lambda_n>0,
\]
because \(c>0\), \(n\ge1\), and \(\lambda_n\ge1\).  Thus division by \(s_n\) is legal.  The exposure realization in \(\mathrm{def:cgd\text{-}assignment}\), arbitrary neighborhood interference, and \(q_{n,i}=\mathbb P_U(E_{n,i,k})>0\) give
\[
 \begin{aligned}
 \mathbb E_UF_n
 &=\sum_{i\in\mathcal I_n}\mathbb E_U\left[
 a_{n,i}\frac{\mathbf 1(E_{n,i,1})}{q_{n,i}}
 -b_{n,i}\frac{\mathbf 1(E_{n,i,0})}{q_{n,i}}
 \right]\\
 &=\sum_{i\in\mathcal I_n}(a_{n,i}-b_{n,i})
 =n\tau_n.
 \end{aligned}
\]
This proves, rather than defines, the equality between the expectation of the observed-data estimator and the fixed causal contrast.

By \(\mathrm{lem:exact\text{-}forward\text{-}doob}\), the variables
\[
 X_{n,s}=M_{n,s}/s_n,\qquad 1\le s\le n,
\]
form a martingale-difference sequence relative to \((\mathcal F_{n,s})_{s=0}^n\), and
\[
 \sum_{s=1}^n\mathbb E_UX_{n,s}^2
 =\frac{\mathbb E_UQ_n}{s_n^2}=1,
 \qquad
 \sum_{s=1}^n\mathbb E_U[X_{n,s}^2\mid\mathcal F_{n,s-1}]
 =\frac{Q_n}{s_n^2}.
\]
The fourth moments are finite by \(\mathrm{lem:prefix\text{-}flip\text{-}moments}\).  Hence the cited Heyde--Brown inequality \(\mathrm{lem:heyde\text{-}brown\text{-}martingale\text{-}fourth\text{-}moment}\) applies and yields
\[
 \begin{aligned}
 d_{\mathrm K}\left(\frac{F_n-\mathbb E_UF_n}{s_n},Z_0\right)
 &\le C\left\{
 \frac{\sum_{s=1}^n\mathbb E_U|M_{n,s}|^4}{s_n^4}
 +\mathbb E_U\left|\frac{Q_n}{s_n^2}-1\right|^2
 \right\}^{1/5}\\
 &=C\left\{
 \frac{\sum_{s=1}^n\mathbb E_U|M_{n,s}|^4}{s_n^4}
 +\frac{\operatorname{Var}_U(Q_n)}{s_n^4}
 \right\}^{1/5}.
 \end{aligned}
\]
The equality uses \(\mathbb E_UQ_n=s_n^2\), again from \(\mathrm{lem:exact\text{-}forward\text{-}doob}\).  Applying the two inequalities in \(\mathrm{lem:prefix\text{-}flip\text{-}moments}\) and the lower bound \(s_n^4\ge c^2n^2\lambda_n^2\) gives, term by term,
\[
 \frac{\sum_{s=1}^n\mathbb E_U|M_{n,s}|^4}{s_n^4}
 \le \frac{C_0B D_n^2}{c^2n\lambda_n},
 \qquad
 \frac{\operatorname{Var}_U(Q_n)}{s_n^4}
 \le \frac{C_0B D_n}{c^2n}.
\]
Since \(F_n=nT_n\), \(\mathbb E_UF_n=n\tau_n\), and \(s_n=n\sigma_n\), the standardized variable on the left is exactly \((T_n-\tau_n)/\sigma_n\).  Absorbing \(C\), \(C_0\), \(B\), and \(c^{-2}\) into \(C_{\mathrm N}(B,c)\), and using that Kolmogorov distance is at most one, proves
\[
 d_{\mathrm K}\left(\frac{T_n-\tau_n}{\sigma_n},Z_0\right)
 \le \min\left\{1,C_{\mathrm N}(B,c)
 \left(\frac{D_n^2}{n\lambda_n}+\frac{D_n}{n}\right)^{1/5}\right\}.
\]

By \(\mathrm{lem:perron\text{-}prefix\text{-}operator}\), \(\lambda_n\ge\sqrt{D_n}\); and \(D_n\ge1\).  Therefore
\[
 \frac{D_n^2}{n\lambda_n}\le\frac{D_n^{3/2}}n,
 \qquad
 \frac{D_n}{n}\le\frac{D_n^{3/2}}n.
\]
Enlarging \(C_{\mathrm N}(B,c)\) by the factor \(2^{1/5}\) proves the second pointwise bound.  Finally, \(D_n\le d\) for every member of \(\mathcal M_n(d;B,c)\).  Taking the two suprema in \(\mathrm{def:uniform\text{-}error}\) therefore gives
\[
 K_n(d;B,c)\le\min\left\{1,C_{\mathrm N}(B,c)
 \left(\frac{d^{3/2}}n\right)^{1/5}\right\},
\]
as claimed.

%%% FIELD statement-graph-adaptive-corollary
For fixed \(B\ge1\), \(0<c\le1\), and \(\alpha\in(0,1)\), let \(((G_n,Y_n))_{n\ge1}\) be any sequence such that \((G_n,Y_n)\in\mathcal M_n(D_n;B,c)\) for every \(n\).  The graph-adaptive condition
\[
 \frac{D_n^2}{n\lambda_n}\longrightarrow0
\]
is sufficient for
\[
 d_{\mathrm K}\left(\frac{T_n-\tau_n}{\sigma_n},Z_0\right)\longrightarrow0,
 \qquad
 \frac{\widehat{\mathrm{VB}}_n}{\mathrm{VB}_n}
 \xrightarrow{\mathbb P_U}1,
 \qquad
 \liminf_{n\to\infty}\mathbb P_U(\tau_n\in C_{n,\alpha})\ge1-\alpha.
\]
For the repeated-star construction at any fixed \(\beta>0\), the same diagnostic converges instead to \(\beta^{3/2}>0\), while the standardized estimator converges to \(\mathcal L_{\beta}\ne\mathcal N(0,1)\).  Thus the repeated-star family is a nonvanishing-diagnostic elbow illustration; the sufficient condition is not asserted to be necessary for each individual graph sequence.

%%% FIELD proof-graph-adaptive-corollary
Set
\[
 r_n=\frac{D_n^2}{n\lambda_n},
 \qquad
 \kappa_n=d_{\mathrm K}\left(\frac{T_n-\tau_n}{\sigma_n},Z_0\right).
\]
By \(\mathrm{lem:perron\text{-}prefix\text{-}operator}\), \(1\le\lambda_n\le D_n\).  Consequently
\[
 \frac{D_n}{n}=r_n\frac{\lambda_n}{D_n}\le r_n,
 \qquad
 \frac{\lambda_n}{n}=r_n\left(\frac{\lambda_n}{D_n}\right)^2\le r_n.
\]
If \(r_n\to0\), \(\mathrm{thm:uniform\text{-}normal\text{-}bound}\), applied pointwise to the member of \(\mathcal M_n(D_n;B,c)\), gives
\[
 \kappa_n\le C_{\mathrm N}(B,c)(2r_n)^{1/5}\longrightarrow0.
\]
This proves the normal limit directly from the primitive design bounds rather than assuming it.

For every fixed \(\eta>0\), \(\mathrm{thm:spectral\text{-}ratio}\) and the second displayed inequality give
\[
 \mathbb P_U\left(
 \left|\frac{\widehat{\mathrm{VB}}_n}{\mathrm{VB}_n}-1\right|>\eta
 \right)
 \le\frac{C_0B\lambda_n}{c^2n\eta^2}
 \le\frac{C_0B r_n}{c^2\eta^2}\longrightarrow0.
\]
Hence \(\widehat{\mathrm{VB}}_n/\mathrm{VB}_n\xrightarrow{\mathbb P_U}1\).  Notice also that \(\mathrm{thm:spectral\text{-}ratio}\) gives \(\mathrm{VB}_n\ge\sigma_n^2>0\), so all ratios used here are well defined.

It remains to prove feasible coverage using the actual error \(\kappa_n\), rather than a worst-case envelope indexed only by \(D_n\).  Put \(z=z_{1-\alpha/2}>0\),
\[
 u_n=\frac{\lambda_n}{n},
 \qquad
 \varepsilon_n=u_n^{1/3},
 \qquad
 \mathcal A_n=\left\{
 \left|\frac{\widehat{\mathrm{VB}}_n}{\mathrm{VB}_n}-1\right|
 \le\varepsilon_n\right\}.
\]
We have \(0<u_n\le1\), because \(1\le\lambda_n\le D_n\le n\), and \(u_n\le r_n\to0\).  Thus eventually \(0<\varepsilon_n<1\).  On \(\mathcal A_n\),
\[
 \widehat{\mathrm{VB}}_n
 \ge(1-\varepsilon_n)\mathrm{VB}_n
 \ge(1-\varepsilon_n)\sigma_n^2.
\]
If \(\widehat{\mathrm{VB}}_n=0\), \(\mathrm{def:wald\text{-}interval}\) returns \(\mathbb R\), so noncoverage is impossible.  Otherwise, with \(X_n=(T_n-\tau_n)/\sigma_n\), the preceding inequality implies
\[
 \mathbb P_U(\tau_n\notin C_{n,\alpha})
 \le \mathbb P_U\{|X_n|>z\sqrt{1-\varepsilon_n}\}
 +\mathbb P_U(\mathcal A_n^{\mathsf c}).
\]
The definition of \(\kappa_n\), applied to the two tails, yields
\[
 \mathbb P_U\{|X_n|>x\}\le2\{1-\Phi(x)\}+2\kappa_n
 \quad(x\ge0).
\]
Because \(2\{1-\Phi(z)\}=\alpha\), the standard normal density is bounded by \((2\pi)^{-1/2}\), and \(1-\sqrt{1-\varepsilon_n}\le\varepsilon_n\), it follows that
\[
 2\{1-\Phi(z\sqrt{1-\varepsilon_n})\}
 \le\alpha+\frac{2z}{\sqrt{2\pi}}\varepsilon_n.
\]
Finally, \(\mathrm{thm:spectral\text{-}ratio}\) with \(\eta=\varepsilon_n\) gives
\[
 \mathbb P_U(\mathcal A_n^{\mathsf c})
 \le\frac{C_0B u_n}{c^2\varepsilon_n^2}
 =\frac{C_0B}{c^2}u_n^{1/3}.
\]
Combining the last four displays,
\[
 \mathbb P_U(\tau_n\notin C_{n,\alpha})
 \le\alpha+2\kappa_n+
 \left(\frac{2z}{\sqrt{2\pi}}+\frac{C_0B}{c^2}\right)u_n^{1/3}.
\]
Both remainder terms tend to zero, proving the claimed conservative lower coverage limit.

For the repeated-star construction of \(\mathrm{def:repeated\text{-}star\text{-}witness}\), direct diagonalization of a star with \(k_n\) leaves gives
\[
 D_n=k_n+1,
 \qquad
 \lambda_n=1+\sqrt{k_n}.
\]
Since \(k_n\sim\beta n^{2/3}\),
\[
 \frac{D_n^2}{n\lambda_n}longrightarrow\beta^{3/2}.
\]
The nonnormal convergence to \(\mathcal L_\beta\) is exactly \(\mathrm{thm:repeated\text{-}star\text{-}boundary}\).  This verifies the stated elbow illustration, while making no converse claim for arbitrary individual graph sequences.
